\documentclass[11pt]{article}

\usepackage[margin=1.1in]{geometry}
\usepackage{amsmath,amssymb,amsthm}
\usepackage{stmaryrd}
\usepackage{mathtools}
\usepackage{tikz}
\usetikzlibrary{arrows.meta,positioning,fit,backgrounds,patterns,calc,decorations.pathmorphing}
\usepackage{listings}
\usepackage{xcolor}
\usepackage{booktabs}
\usepackage{array}
\usepackage[colorlinks=true,linkcolor=black,citecolor=black,urlcolor=blue]{hyperref}
\usepackage[expansion=false]{microtype}

\theoremstyle{plain}
\newtheorem{theorem}{Theorem}
\newtheorem{proposition}[theorem]{Proposition}
\newtheorem{lemma}[theorem]{Lemma}
\newtheorem{corollary}[theorem]{Corollary}
\newtheorem{conjecture}[theorem]{Conjecture}
\theoremstyle{definition}
\newtheorem{definition}[theorem]{Definition}
\newtheorem{remark}[theorem]{Remark}
\newtheorem{observation}[theorem]{Observation}
\newtheorem{question}{Question}

% --- typography: the rho calculus is NEVER written with the Greek letter ---
\newcommand{\rhoc}{\textsf{rho}}

\lstdefinestyle{rho}{
  basicstyle=\ttfamily\small,
  keywordstyle=\bfseries,
  morekeywords={for,new,in,def,where,contract,match,if,else,Nil},
  columns=fullflexible,
  showstringspaces=false,
  breaklines=true,
  frame=none,
  xleftmargin=1.5em,
  comment=[l]{//},
  commentstyle=\itshape\color{black!55},
}
\lstset{style=rho}

\title{\bfseries The Gap\\[.35em]
\large Deception as a bootstrap mechanism for grounded symbol systems}
\author{L.\ G.\ Meredith\\ \small F1R3FLY.io}
\date{September 8, 2026}

\begin{document}
\maketitle

\begin{abstract}
An earlier note established a negative result: getting the distribution of a symbol system is not
getting its correspondence to the world, and the correspondence is carried in the community of
speakers rather than in the symbol system. Linear~A and Rongorongo supply the corpus without the
community and resist interpretation; whale song supplies the community without a Rosetta stone and is
yielding. That note also observed that the bootstrap problem must be soluble, since humans and whales
solved it independently, and isolated three features common to both solutions. It did not say how
either solution was reached. This note records a candidate mechanism, due to Michael Benedikt: that
language originates in lying. A reflexive call and the condition triggering it are welded together and
require no interpretation. Prise them apart --- as a hominid does the first time she gives the snake
call knowing there is no snake, and takes the larger share while the band is in the trees --- and
semantics becomes both possible and necessary. We set the parable inside the mortal scientist
framework, note that its central move is a re-typing of a channel rather than any new construct, and
argue for a thesis: meaning is not a property of the sign but the theory receivers are forced to build
once senders are capable of lying. Two consequences are stated as targets. First, that a corpus
generated by strategic senders is a mixture over sender policies whose weights the corpus does not
identify, which would upgrade the earlier note's central observation from an observation to a theorem.
Second, that grounding is therefore a \emph{drive} rather than a \emph{read} --- available only to an
agent that can emit into a community and observe the response. The relevant ethological literature is
better than anecdotal and is recorded here, as is the fact that the deception route is not the only
one.
\end{abstract}

\section{Where this picks up}
\label{sec:pickup}

The earlier note established the negative result and its two-sided evidence.

\paragraph{Linear~A and Rongorongo.} A corpus large enough to train a transformer to emit well-formed
utterances, and no community of speakers left. Decades of resistance to interpretation. The
distribution survived; the correspondence did not.

\paragraph{Whale song.} No Rosetta stone and no cognates, and yet real progress, because the population
is alive and its behaviour can be observed. Rongorongo minus the extinction.

The conclusion drawn there was that the correspondence between a symbol system and the world is carried
in the community of speakers. Decipherment is transport into an anchor, or direct observation of the
using population; neither route is distributional. The note further recorded that the bootstrap problem
\emph{has} to be soluble, and isolated three features common to the human and cetacean solutions: the
solver was a population, the interface was not chosen, and there was an external criterion supplying the
fitness signal that selects among relabellings.

What that note did not supply was a mechanism. This one records a candidate.

\section{Benedikt's argument}
\label{sec:benedikt}

Michael Benedikt was ACSA Distinguished Professor of Architecture at the University of Texas at Austin,
held the Hal Box Chair in Urbanism, and wrote \emph{Deconstructing the Kimbell} (1991) among much else.
He died in August 2025. The argument below is recalled from conversation and should be cited as personal
communication; a first search surfaced no published version.

Benedikt argues that language originates in lying.

\paragraph{The setting.} A band of early hominids on the plain, not far from the trees, dividing food:
fruit, or the spoils of a kill.

\paragraph{The prior state.} Many animals make reflexive noises grounded in function --- the sounds of
breathing, howling in pain, the alarm call. These are triggered by conditions. They are not meaningful
in the way that language is meaningful. Benedikt stresses this, and it bears stressing again, because
the whole argument turns on what the prior state lacks.

\paragraph{The event.} Eve notices movement in the grass. The conditions land on her nervous system in
just such a way as to trigger the snake warning call. She calls; the band scatters for the trees. In
that split second Eve determines that it was not a snake --- it was the wind --- and at the same moment
notices that she now stands in an advantageous position with respect to the division of the food.

What she sees is \emph{the gap}: the distance between the reflexively triggered sound and the behaviour
of the band. And she sees that the gap admits exploitation.

\paragraph{The follow-on.} If Eve does this regularly, others may notice the correlation between the
snake call and the division of the food. The band learns Eve's trick from Eve's behaviour. Language is a
virus \cite{burroughs}.

\paragraph{The reading.} It is the gap that gives the opening for meaning. A reflexive call and its
associated condition are welded together and nothing there needs interpreting. Prise them apart and
semantics becomes possible. This is the sense in which lying is the knowledge of good and evil offered
at the centre of the garden.

\section{The empirical situation is better than anecdotal}
\label{sec:empirical}

The parable has data behind it, and the data sharpen it in a direction the parable does not anticipate.

\paragraph{Tufted capuchins.} Wheeler \cite{wheeler2009} reports that wild tufted capuchins give
predator calls in non-predatory contexts where the caller stands to gain. False alarms are given by
subordinates more than by dominants, rise with the contestability of the food and with its scarcity,
and are given from positions in which the caller would gain if the others reacted. Listeners leave the
patch; the caller takes the food. This is Eve's scenario, observed.

\paragraph{Counterdeception.} Listeners are more likely to ignore the same call type when it is produced
during competitive feeding than in other contexts, and the difference is carried primarily by a drop in
escape responses rather than in vigilance \cite{wheeler2010,wheeler2013}. The receivers have become
context-sensitive about a signal that used to require no context.

\paragraph{Fork-tailed drongos.} Drongos steal food using \emph{mimicked} alarm calls of other species,
and vary which call they mimic \cite{flower2011,flower2014}. The stated reason is that a deceptive
signal loses power as its frequency rises relative to honest signals, so varying the signal is what
sustains the deception.

\paragraph{Earlier work.} Munn's kleptoparasitic shrike-tanagers and antshrikes \cite{munn1986}, whose
status as alarm rather than aggressive signals has since been questioned; and arctic foxes calling cubs
off food.

\begin{observation}[Deception grows the vocabulary]
\label{obs:vocab}
Under frequency-dependent discounting, deception does not merely exploit an existing repertoire. It
exerts pressure to expand the repertoire, because a fresh signal starts at the honest response rate.
\end{observation}

Observation~\ref{obs:vocab} is a mechanism for vocabulary growth, not only for the birth of meaning, and
it is the part of the picture that the parable alone would not have produced.

\section{The framework translation}
\label{sec:translation}

\subsection{Three stages}

\paragraph{Stage 0: the reflex.} The call is a rewrite rule, not a symbol. In \rhoc{} it is a
sensor-driven emission, and the receiver's response is another rewrite. The correspondence to the world
is a fact about the wiring. Neither party needs a hypothesis. In the ladder's terms both sides are
non-modal: this is perception, reading structure the world has already computed into legible form.

\paragraph{Stage 1: the decoupling.} Eve's innovation is to make the call channel writable by her policy
independently of the sensor. In the channel classification of the scientist note --- channels enabling
internal reductions, channels enabling interaction with the environment, and channels doing both ---
this is a \emph{re-typing}. An emission driven only from the sensor side becomes driven from the policy
side as well. Nothing is added to the calculus. The gap is a change in which redexes can fire on a name.

\paragraph{Stage 2: semantics as the receiver's forced expense.} Once emission is decoupled, ``call
heard'' no longer entails ``snake present''. The receiver's predicate cannot remain non-modal. It must
be indexed by context, and the context is a hypothesis about the \emph{sender} rather than about the
world. This is exactly the machinery already present: context-labelled modalities, and an environment
whose most interesting contents are other scientists.

\begin{observation}[Thesis]
\label{obs:thesis}
Meaning is not a property of the sign. Meaning is the theory that receivers are forced to build once
senders are capable of lying. Semantics is the receiver's bill.
\end{observation}

\subsection{Statements to attempt}

\begin{proposition}[No semantics before the lie]
\label{prop:A}
In a population where every emission is sensor-driven, the receiver's optimal policy is a non-modal
predicate on the signal: the modal rungs of the ladder buy nothing.
\end{proposition}

The corollary is that an alarm call in a wholly honest population is not a symbol --- and the framework
can say so with a purchase criterion rather than with a definition, which is the difference between a
result and a stipulation.

\begin{proposition}[The lie forces the modal rung]
\label{prop:B}
With a nonzero deception rate the non-modal predicate is strictly dominated, and the receiver's best
available response requires context-indexing. The transition has a threshold in the deception rate.
\end{proposition}

The threshold should be computable in the currency of the game note's ladder returns, and
Propositions~\ref{prop:A} and \ref{prop:B} together are the pair a ladder computation could actually
decide.

\begin{proposition}[The deception equilibrium]
\label{prop:C}
The liar's yield is proportional to the response rate, which declines in the population deception rate;
hence an interior equilibrium.
\end{proposition}

The empirical predictions are already made and already met: subordinates lie more, lying rises with
contestability, listeners discount contextually \cite{wheeler2009,wheeler2010}.

\begin{proposition}[Elaboration under discounting]
\label{prop:D}
At equilibrium there is pressure to diversify the signal. Vocabulary size grows as a function of the
deception pressure and of the cost of manufacturing a new distinguishable name.
\end{proposition}

Freshness by quotation says names are cheap to manufacture, which predicts fast growth --- plausibly too
fast. What limits it is a real question and a good one.

\begin{proposition}[The trick is horizontally transmitted code]
\label{prop:E}
The band learns Eve's trick by watching Eve, not by descending from her: a quoted process copied between
contemporaries. This is the reproduction material with the vertical arrow replaced by a horizontal one.
\end{proposition}

Burroughs' virus is not a metaphor under Proposition~\ref{prop:E}. It is an $R_{0}$, and the thing it
spreads degrades the resource it spreads on.

\begin{proposition}[Strategic use makes the corpus unidentifiable]
\label{prop:F}
If emission is policy-driven, the joint distribution over signals and world-states is a mixture over
sender policies whose weights are facts about the population. A corpus gives the marginal, and the
marginal does not identify the mixture.
\end{proposition}

Proposition~\ref{prop:F} is the statement with reach outside the framework, because it says that no
corpus of any size recovers the correspondence, and says it as a theorem about strategic signalling
rather than as a complaint about data scarcity. If it goes through, the earlier note's central
observation is upgraded: the correspondence lives in the community \emph{because} the senders are
strategic. Linear~A resists because its semantics was never in the distribution. Whale song yields
because the senders remain available to be modelled.

\begin{proposition}[Grounding is a drive, not a read]
\label{prop:G}
A transformer over a corpus can only read. Under Proposition~\ref{prop:F} the readable object is the
wrong one, and the identifying information is available only to an agent that can emit into the
community and observe the response --- that is, to one that can perform the intervention separating the
mixture components.
\end{proposition}

Grounding then requires participation rather than perception. That is a sharp and unwelcome claim about
what embodiment has to mean, and its interest is that it is stated in the framework's own pricing
vocabulary rather than imported.

\section{Caveats and honest gaps}
\label{sec:caveats}

\begin{enumerate}
\item \textbf{Deception is not the only bootstrap.} Lewis conventions and Skyrms' signalling games
produce meaning by reinforcement under aligned interests, with no liar anywhere \cite{lewis,skyrms}. The
defensible claim is narrower than ``lying is necessary for signalling'': it is that \emph{lying is what
makes meaning require modelling}, forcing the receiver up from correlation to theory. The two routes
should yield different kinds of symbol system, and characterising the difference is a good problem.

\item \textbf{Deception is parasitic on an honest baseline.} The reflexive layer must exist first;
otherwise there is nothing to exploit. This constrains any engineered bootstrap: build a functionally
grounded reflex layer first, before there is anything for a policy to detach from.

\item \textbf{The ethological alignment should be checked rather than reinvented.} Dawkins and Krebs
\cite{dawkinskrebs} argue that signals are manipulation rather than information transfer, which is
Benedikt's thesis arriving from ethology; and the equilibrium of Proposition~\ref{prop:C} has
presumably been worked out in the costly-signalling literature \cite{msharper} in a different currency.
Import it and pay it in tokens.

\item \textbf{Intentionality is not claimed and should not be.} The capuchin literature is careful about
whether false alarms are deliberate, to the point of testing whether adrenocortical stress explains the
calling. The framework does not need deliberateness. It needs only that emission be policy-driven rather
than sensor-driven, which is a structural condition on channels and is exactly what
Propositions~\ref{prop:A} and \ref{prop:B} are stated in terms of.

\item \textbf{The gap-as-opening reading does philosophical work the formal apparatus may not reach.}
Keep the Genesis line as a closing remark rather than as a premise.
\end{enumerate}

\section{Open questions}
\label{sec:open}

\begin{question}
Does the deception rate play the role in Propositions~\ref{prop:A}--\ref{prop:B} that the environment
plays in the compose note's inversion --- so that whether \emph{meaning} is affordable is a property of
the environment rather than of the learner, exactly as whether \emph{truth} is affordable turned out to
be?
\end{question}

\begin{question}
Can Proposition~\ref{prop:F} be given as a genuine identifiability theorem, with the mixture stated over
a space of sender policies and the failure of identifiability exhibited rather than argued?
\end{question}

\begin{question}
The earlier note's third common feature --- an external criterion supplying the fitness signal --- is
here instantiated as the food share. Its first feature --- that the solver was a population --- is here
\emph{explained} rather than observed, since a lie requires someone to lie to. Does the second feature,
that the interface was not chosen, receive a similar treatment?
\end{question}

\begin{question}
Open source game theory is the limit in which deception is impossible, since agents reason over source
rather than over claims. In the three grades of access that is the reflection grade, and deception lives
at perception and interaction. Does deception return under bounded verification, when the cost of
checking a property exceeds the stake --- making obfuscation an economic attack rather than an
informational one?
\end{question}

\begin{question}
In the collapse result for the two open-source bots, the pair are extensionally identical on the
decidable fragment and differ only in their $\bot$-disposition. Deception appears to live entirely in
that complement. Is the framework-native definition available: \emph{deception is engineering the
opponent into $\bot$ and harvesting their default}?
\end{question}

\section{Filing}
\label{sec:filing}

Candidate fifth note in the rho-life series, or a substantial new section of the bootstrap note.
Propositions~\ref{prop:C}, \ref{prop:D} and \ref{prop:F} carry measurable content and an attached
empirical literature, which is more than a section can comfortably hold.

\begin{thebibliography}{99}

\bibitem{burroughs} W.~S.~Burroughs, \emph{The Ticket That Exploded}, Olympia Press, 1962. The
formulation used here is Laurie Anderson's, from \emph{Home of the Brave}, 1986.

\bibitem{dawkinskrebs} R.~Dawkins and J.~R.~Krebs, Animal signals: information or manipulation?, in
\emph{Behavioural Ecology: An Evolutionary Approach}, Blackwell, 1978.

\bibitem{flower2011} T.~Flower, Fork-tailed drongos use deceptive mimicked alarm calls to steal food,
\emph{Proceedings of the Royal Society~B} \textbf{278} (2011), 1548--1555.

\bibitem{flower2014} T.~P.~Flower, M.~Gribble and A.~R.~Ridley, Deception by flexible alarm mimicry in
an African bird, \emph{Science} \textbf{344} (2014), 513--516.

\bibitem{lewis} D.~Lewis, \emph{Convention: A Philosophical Study}, Harvard University Press, 1969.

\bibitem{msharper} J.~Maynard Smith and D.~Harper, \emph{Animal Signals}, Oxford University Press,
2003.

\bibitem{munn1986} C.~A.~Munn, Birds that ``cry wolf'', \emph{Nature} \textbf{319} (1986), 143--145.

\bibitem{skyrms} B.~Skyrms, \emph{Signals: Evolution, Learning, and Information}, Oxford University
Press, 2010.

\bibitem{wheeler2009} B.~C.~Wheeler, Monkeys crying wolf? Tufted capuchin monkeys use anti-predator
calls to usurp resources from conspecifics, \emph{Proceedings of the Royal Society~B} \textbf{276}
(2009), 3013--3018.

\bibitem{wheeler2010} B.~C.~Wheeler, Decrease in alarm call response among tufted capuchin monkeys in
competitive feeding contexts: possible evidence for counterdeception, \emph{International Journal of
Primatology} \textbf{31} (2010).

\bibitem{wheeler2013} B.~C.~Wheeler and K.~Hammerschmidt, Proximate factors underpinning receiver
responses to deceptive false alarm calls in wild tufted capuchin monkeys: is it counterdeception?,
\emph{American Journal of Primatology} \textbf{75} (2013), 715--725.

\end{thebibliography}

\end{document}
